\documentclass[10pt,a4paper]{amsart}
\usepackage[margin=19mm]{geometry}
\usepackage{amsmath,amssymb,mathtools}
\usepackage[scr=rsfs,cal=euler]{mathalfa}
\usepackage{xfrac}
\usepackage[unicode,hidelinks,bookmarks=false]{hyperref}
\newcommand{\ie}{i.e.\ }
\newcommand{\cf}{cf.\ }
\newcommand{\resp}{respectively\ }
\newcommand{\Subsection}{\subsection}
\providecommand{\nobreakdash}{\nobreak\hbox{-}}
\setlength{\emergencystretch}{2em}
\multlinegap0pt
\usepackage{amssymb}
\newtheorem{theorem}{Theorem}
\title{A formula of counting divisors in integers rings: a generalization of the divisor function $d_0(n)$}
\author{\'Angel Mart\'inez-Avelar, Mario Pineda-Ruelas}
\date{Source: arXiv:2605.19177v1 (CC BY 4.0)}
\begin{document}
\maketitle

\noindent\emph{Source and licence.} A formula of counting divisors in integers rings: a generalization of the divisor function $d_0(n)$. Version arXiv:2605.19177v1 (CC BY 4.0), distributed by arXiv under the Creative Commons Attribution 4.0 International licence, http://creativecommons.org/licenses/by/4.0/, as recorded in the arXiv licence metadata for this exact version and shown on the abstract page https://arxiv.org/abs/2605.19177v1. Full licence notice: \texttt{licenses/arxiv-2605.19177-CC-BY-4.0.txt}.\par\smallskip

\noindent\textit{Editorial scope.} Counted unit: the article's theorem teoprinc (source0.tex lines 130--134). The article's complete original proof of it is delivered with it (source0.tex lines 136--160, one \texttt{proof} environment of 25 lines). No mathematics is written, repaired or re-derived: the statement and the proof are the article's own lines, in the article's own notation, and no retained line is substituted or re-flowed.  No further source-local context is needed: every cross-reference of every retained line resolves inside the two delivered spans.  Nothing was omitted from any retained span, so the package carries no omission record.  No retained line contains a citation, so the wrapper carries no bibliography and no bibliography entry.  No retained line contains a figure environment, an image-inclusion command or drawing code, so no image is delivered and the package declares no asset dependency.  Editorial formatting only, in addition: an AMS \texttt{amsart} preamble that restates the article's own declarations and macros used by the retained lines, namely the environments \texttt{theorem} on separate counters, together with the standard packages those lines need.  The retained environments print in this document as Theorem 1 in order of appearance, whereas the article prints the same items with the numbering of its own full document.  The complete native source of the article as distributed in the arXiv e-print for this exact version is bundled unchanged as \texttt{sources/200899/source0.tex}.
\begin{theorem}\label{teoprinc}
Let $K$ be a number field, $\mathcal{O}_K$ its integers ring and $I \subseteq \mathcal{O}_K$ an ideal. Suppose that $Cl(K)$ is cyclic of order $h$. Then there exist $\gamma \in \mathcal{O}_K$, distinct ideals $T_1, \dots, T_{m}$ and $c_1,\ldots,c_m\in\mathbb{N}$ such that $c_1+\ldots+c_m\leq h-1$ satisfying
\[ I = (\gamma) T_1^{c_1} \cdots T_{m}^{c_m}, \]
where for each $i$, $T_i$ is a non-principal prime ideal or $T_i=\mathcal{O}_K$.
\end{theorem}

\begin{proof}
   If $I=\mathcal{O}_K$ or $I=\{0\}$, then the result is trivial. Suppose that $Cl(K) \cong \mathbb{Z}/h\mathbb{Z}$, and $I=(\alpha,\beta)$ is a non-zero proper ideal of $\mathcal{O}_K$. Consider the following prime ideal factorizations:
\[ (\alpha) = P_1^{a_1} \cdots P_g^{a_g}, \]
\[ (\beta) = Q_1^{b_1} \cdots Q_r^{b_r}. \]
Suppose that $(\alpha)$ and $(\beta)$ have no common principal ideal divisor other than the trivial one. Since $I\neq \mathcal{O}_K$, there exist $1\leq i\leq g$ and $1\leq j\leq r$ such that $P_i=Q_j$. Without loss of generality, assume such equalities occur in the first $m$ factors, i.e., $P_i=Q_i$ for $i=1,\ldots, m$. For each $1\leq i\leq m$, we let $T_i=P_i=Q_i$ and $c_i=\min\{a_i,b_i\}$. Then

\begin{equation*}
\begin{split}
I &= (\alpha,\beta) = (\alpha) + (\beta) \\
  &= T_1^{c_1} \cdots T_m^{c_m} (P_1^{\alpha_1} \cdots P_m^{\alpha_m} P_{m+1}^{a_{m+1}} \cdots P_g^{a_g} + Q_1^{\beta_1} \cdots Q_m^{\beta_m} Q_{m+1}^{b_{m+1}} \cdots Q_r^{b_r}),
\end{split}
\end{equation*}
where each $T_i$ is non-principal, $T_i\neq T_j$ if $i\neq j$ and $c_i+\alpha_i=a_i$, $c_i+\beta_i=b_i$ for $i=1,\ldots, m$. Since there are no more common factors in the sum, we have

\[ I=(1)T_1^{c_1}\cdots T_m^{c_m}. \] 
If $c_1+\ldots+c_m\geq h$, then $[T_1]^{c_1}\cdots [T_m]^{c_m}$ is a sequence of $h$ or more elements in $Cl(K)$. Given that $\mathcal{D}(Cl(K))=h$, there exist $i_1,\ldots,i_s$ such that $[T_{i_1}]\cdots[T_{i_s}]=1$, which is not possible because there are no non-trivial principal ideal divisors. Therefore $c_1+\ldots+c_m\leq h-1$. Suppose there exists a non-trivial common principal ideal divisor $(\gamma)$ of $(\alpha)$ and $(\beta)$, and we choose $(\gamma)$ to be maximal. Then 

\[ I=(\alpha,\beta)=(\gamma)(a,b), \]
where $(\gamma a)=(\alpha)$, $(\gamma b)=(\beta)$ and the ideals $(a)$, $(b)$ have no common principal ideal divisors by the maximality of $(\gamma)$. From the previous case, we have 

\[ (a,b)=(1)T_1^{c_1}\cdots T_m^{c_m}, \]
where $c_1+\ldots+c_m\leq h-1$ and $T_i$ is a non-principal prime ideal for all $i$. Therefore 

\[ I=(\gamma)T_1^{c_1}\cdots T_m^{c_m}. \]
\end{proof}

\end{document}
